\documentclass[12pt,letterpaper]{article}
\usepackage{graphicx,textcomp}
\usepackage{natbib}
\usepackage{setspace}
\usepackage{fullpage}
\usepackage{color}
\usepackage[reqno]{amsmath}
\usepackage{amsthm}
\usepackage{fancyvrb}
\usepackage{amssymb,enumerate}
\usepackage[all]{xy}
\usepackage{endnotes}
\usepackage{lscape}
\newtheorem{com}{Comment}
\usepackage{float}
\usepackage{hyperref}
\newtheorem{lem} {Lemma}
\newtheorem{prop}{Proposition}
\newtheorem{thm}{Theorem}
\newtheorem{defn}{Definition}
\newtheorem{cor}{Corollary}
\newtheorem{obs}{Observation}
\usepackage[compact]{titlesec}
\usepackage{dcolumn}
\usepackage{tikz}
\usetikzlibrary{arrows}
\usepackage{multirow}
\usepackage{xcolor}
\newcolumntype{.}{D{.}{.}{-1}}
\newcolumntype{d}[1]{D{.}{.}{#1}}
\definecolor{light-gray}{gray}{0.65}
\usepackage{url}
\usepackage{listings}
\usepackage{color}

\definecolor{codegreen}{rgb}{0,0.6,0}
\definecolor{codegray}{rgb}{0.5,0.5,0.5}
\definecolor{codepurple}{rgb}{0.58,0,0.82}
\definecolor{backcolour}{rgb}{0.95,0.95,0.92}

\lstdefinestyle{mystyle}{
	backgroundcolor=\color{backcolour},   
	commentstyle=\color{codegreen},
	keywordstyle=\color{magenta},
	numberstyle=\tiny\color{codegray},
	stringstyle=\color{codepurple},
	basicstyle=\footnotesize,
	breakatwhitespace=false,         
	breaklines=true,                 
	captionpos=b,                    
	keepspaces=true,                 
	numbers=left,                    
	numbersep=5pt,                  
	showspaces=false,                
	showstringspaces=false,
	showtabs=false,                  
	tabsize=2
}
\lstset{style=mystyle}
\newcommand{\Sref}[1]{Section~\ref{#1}}
\newtheorem{hyp}{Hypothesis}

\title{Problem Set 1 \\
\large Applied Stats/Quant Methods }
\date{Due: October 9, 2025}
\author{Matilda Tomatis}

\begin{document}
	\maketitle
	
	\section*{Instructions}
	\begin{itemize}
	\item Please show your work! You may lose points by simply writing in the answer. If the problem requires you to execute commands in \texttt{R}, please include the code you used to get your answers. Please also include the \texttt{.R} file that contains your code. If you are not sure if work needs to be shown for a particular problem, please ask.
\item Your homework should be submitted electronically on GitHub.
\item This problem set is due before 23:59 on Thursday October 9, 2025. No late assignments will be accepted.
	\end{itemize}
	
	\vspace{1cm}
	\section*{Question 1: Education}

A school counselor was curious about the average of IQ of the students in her school and took a random sample of 25 students' IQ scores. The following is the data set:\\
\vspace{.5cm}

\lstinputlisting[language=R, firstline=35, lastline=35]{PS01_Tomatis.R}  

\vspace{1cm}

\begin{enumerate}
	\item Find a 90\% confidence interval for the average student IQ in the school.\\
	
	\item Next, the school counselor was curious  whether  the average student IQ in her school is higher than the average IQ score (100) among all the schools in the country.\\ 
	
	\noindent Using the same sample, conduct the appropriate hypothesis test with $\alpha=0.05$.
\end{enumerate}
\vspace{1cm}
\noindent\textbf{Answer 1}\\

\noindent There are different steps needed in order to properly construct a 90\% confidence interval. Firstly we will need to find the sample parameter in question, hence we will calculate the sample mean out of the students' IQs.\\
\vspace{.5cm}
\lstinputlisting[language=R, firstline=40, lastline=41]{PS01_Tomatis.R}  
\vspace{.5cm}
\noindent Subsequently we calculate standard deviation and standard error from the sample. \\
\vspace{.5cm}
\lstinputlisting[language=R, firstline=45, lastline=48]{PS01_Tomatis.R} 
\vspace{.5cm}
\noindent Moreover, we compute the t-value related to a 90\% C.I. This computation is done using the t-distribution as the sample is smaller than 30. With n = 25 the degrees of freedom are 24, the corresponding distribution is used to compute the confidence interval. \\
\vspace{.5cm}
\lstinputlisting[language=R, firstline=53, lastline=54]{PS01_Tomatis.R}  
\vspace{.5cm}
\noindent Once all this parameters are acquired, the confidence interval can be constructed by summing (upper bound) and subtracting (lower bound) the product of t-value and SE from the t90 value. The interval found can be interpreted as follows: out of 100 sample average class IQs computed, 90 of them would fall between the lower and upper bound of the CI.\\
\vspace{.5cm}
\lstinputlisting[language=R, firstline=59, lastline=61]{PS01_Tomatis.R}  
\begin{verbatim}
	CI_IQ
	[1]  93.95993 102.92007
\end{verbatim}
\vspace{.5cm}
\noindent\textbf{Answer 2}\\


\vspace{0.5cm}
\noindent We shall start by clearly stating the hyphothesis under examination: \\
\vspace{.5cm}
\begin{enumerate}
	\item Ho states: the average IQ of the class under examination is smaller than or equal to 100\\
	\item Ha states: the average IQ of the class is higher than 100.\\ 
\end{enumerate}
\vspace{.5cm}
\noindent In order to test our hyphothesis we must compute the test statistic, then calculate the a p-value associated with it. As mentioned above, the appropriate distribution for the calculation of p-value is a t-distribution with 24 degrees of freedom. \\
\vspace{.5cm}
\lstinputlisting[language=R, firstline=68, lastline=71]{PS01_Tomatis.R} 
\vspace{.5cm}
\noindent So as to evaluate the result we must consider wheter the p-value is bigger than the set $\alpha=0.05$ \\
\vspace{.5cm}
\lstinputlisting[language=R, firstline=75, lastline=77]{PS01_Tomatis.R} 
\begin{verbatim}
print(result)
[1] "p-value: 0.721538341962394 - Ho cannot be rejected"
\end{verbatim}
\textit{The result of our test suggests that $H_0$ cannot be rejected at the 5\% significance level.} We do not have evidence supporting the claim that the class average IQ is higher than the average IQ score among all schools in the country.
\\ 

\newpage

	\section*{Question 2: Political Economy}

\noindent Researchers are curious about what affects the amount of money communities spend on addressing homelessness. The following variables constitute our data set about social welfare expenditures in the USA. \\
\vspace{.5cm}


\begin{tabular}{r|l}
	\texttt{State} &\emph{50 states in US} \\
	\texttt{Y} & \emph{per capita expenditure on shelters/housing assistance in state}\\
	\texttt{X1} &\emph{per capita personal income in state} \\
	\texttt{X2} &  \emph{Number of residents per 100,000 that are "financially insecure" in state}\\
	\texttt{X3} &  \emph{Number of people per thousand residing in urban areas in state} \\
	\texttt{Region} &  \emph{1=Northeast, 2= North Central, 3= South, 4=West} \\
\end{tabular}


\vspace{.5cm}
\noindent Explore the \texttt{expenditure} data set and import data into \texttt{R}. 
\vspace{.5cm}
\begin{itemize}

\item
Please plot the relationships among \emph{Y}, \emph{X1}, \emph{X2}, and \emph{X3}? What are the correlations among them (you just need to describe the graph and the relationships among them)?
\vspace{.5cm}
\item
Please plot the relationship between \emph{Y} and \emph{Region}? On average, which region has the highest per capita expenditure on housing assistance?
\vspace{.5cm}
\item
Please plot the relationship between \emph{Y} and \emph{X1}? Describe this graph and the relationship. Reproduce the above graph including one more variable \emph{Region} and display different regions with different types of symbols and colors.
\end{itemize}
\vspace{1cm}
\noindent\textbf{Answer 1}\\

\noindent Below are the scatterplots obtained for each pair of Y and X variables in our dataset. Each plot includes the annotated correlation of the two variables in question. The measure was chosen as it allows for comparison across different samples and variables.\\
\vspace{.5cm}
\lstinputlisting[language=R, firstline=98, lastline=107]{PS01_Tomatis.R}  
\vspace{.5cm}
\begin{figure}[h!]\centering
	\caption{\footnotesize Scatterplot of all Y and X variable pairs}
	\label{fig:combined_plots}
	\includegraphics[width=1.10\textwidth]{combined_plots.pdf}
\end{figure}


\vspace{0.5cm}
\noindent Observing the graphs a weak-to-moderate linear tendency is observed in the relationship between the variables, though the graphs do present differences in the intensity of it. A common observation can be made: all linear relationships between the variables are positive. \\
\noindent Notably, the pairs X1 and X2, as well as X2 and X3, have particularly low correlation coefficients (0.2056 and 0.2210), suggesting little shared linear change in those pairs. Thus both per-capita income and the number of financially insecure residents, as well as the latter and the number of people residing in urban areas, seem to vary independently in linear terms. This is evident in the scatterplots as well: a clear linear tendency is not observed in the points. \\
\noindent The other pairwise relationships have overall higher correlations, peaking at 0.5953 for the X1-X3 relationship. This suggest a considerable shared linear change in per capital income and the number of resident in urban areas in a given state. Quite strong is also the correlation between Y and X1 (0.5317), suggesting more affluent states also have higher expenditure in shelters/housing assitance. \\
\noindent The correlations between Y and X2 and between Y and X3, stand on similar moderate grounds (0.4483 and 0.4637). It's worth mentioning that in the first of the two scatterplot, a higher concentration of points with lower values of X2 and higher values of Y is noticeable, suggesting a possible U-Shaped relationship between per capita expenditure in shelters/housing assistance and the number of financially insecure residents. It is be noted, however, that such relation cannote be verified with the correlation value as it only measures linear association. \\
\vspace{1cm}

\noindent\textbf{Answer 2} \\

\noindent In order to plot the levels of expenditure in each region, we first recoded the original numeric variable into a factor. We assigned a name to each of the four numberic codes, thus creating a factore whose levels are the region names. Rather than continous, the newly modified variable is categorical. \\

\lstinputlisting[language=R, firstline=168, lastline=172]{PS01_Tomatis.R}  
\vspace{.5cm}
\noindent We proceeded calculating the average expenditure for each of the regions. Afterwards, plotting the relationship between the Y and Region variable.  \\
\lstinputlisting[language=R, firstline=178, lastline=178]{PS01_Tomatis.R}  
\begin{verbatim}
Northeast North Central         South 
79.44444      83.91667      69.18750 
West 
88.30769 
\end{verbatim}
\lstinputlisting[language=R, firstline=186, lastline=191]{PS01_Tomatis.R}  
\vspace{.5cm}
\begin{figure}[H]\centering
	\caption{\footnotesize Expenditure by region}
	\label{fig:Y_Region}
	\includegraphics[width=.75\textwidth]{Y_Region.pdf}
\end{figure}
\noindent The region which, on average, has a higher per capita expenditure on shelter/housing assistance is the West. It is worth mentioning that the sample observations from this area are, however, very spread out across different values of Y, with a standard deviation of 24.06, which is also the highest among all regions.\\

\vspace{1cm}

\noindent\textbf{Answer 3}\\

\vspace{0.5cm}
\noindent The graph was reported and commented above, in the answer 1 of question 2. We will thus proceed with the creation of the upgraded version as requested.\\
\lstinputlisting[language=R, firstline=199, lastline=204]{PS01_Tomatis.R}  
\vspace{.5cm}

\begin{figure}[H]\centering
	\caption{\footnotesize Scatterplot Y and X, by Region}
	\label{fig:Y_X1_Region}
	\includegraphics[width=.75\textwidth]{Y_X1_R.pdf}
\end{figure}

\vspace{.5cm}
\noindent To better comment on how the relationship between X and Y varies across region we calculated the correlation for each of the 4 areas. 
\lstinputlisting[language=R, firstline=212, lastline=214]{PS01_Tomatis.R}
\begin{verbatim}
  Fact_Region   correlation
<fct>               <dbl>
1 Northeast           0.802
2 North Central       0.184
3 South               0.556
4 West                0.305
\end{verbatim}

\vspace{.5cm}
\noindent The magnitude of the linear relationship between Y and X1 varies greatly across Regions, even though all values are positive. \\
\noindent The Northeast region presents as the one in which income per capita is more strongly related to expenditure on shelters or housing assistance, with a strong correlation equal to 0.802.\\
\noindent On the other hand the North Central region appears to display a rather weak linear relationship between Y and X1; characteristic easily observable in the scatterplot: all the green triangular point belonging to the NC region are nested in the center of the graph. 
\end{document}
